\chapter{Resultados experimentales}

Los resultados proceden de ejecuciones de CoppeliaSim lanzadas
desde Python. Cada ejecución exporta señales a CSV y JSON, y de ahí salen las
tablas y figuras. Las capturas renderizadas desde CoppeliaSim complementan las
gráficas y vinculan cada métrica con la escena simulada.

\section{Validación funcional de la escena}

La validación principal de \path{Sim_T2_Phase1_Basic.ttt} terminó con
\texttt{passed=true}. Los indicadores más relevantes fueron:

\begin{longtable}{P{0.42\textwidth}P{0.34\textwidth}}
\toprule
\textbf{Indicador} & \textbf{Resultado}\\
\midrule
\endhead
Estado final de tarea & \texttt{CHARGING}\\
Modo final de movimiento & \texttt{ARRIVED\_TARGET}\\
Tareas completadas & 4/4\\
Distancia caminada por Bill & \SI{7.843}{m}\\
Batería final R1 & 66,76\%\\
Sensores activos & 16\\
Rayos de sensor visualizados & 16\\
Landmarks/celdas SLAM finales & 24 en validación principal\\
Actualizaciones SLAM & 8532\\
Error final de pose & \SI{0.030}{m}\\
Error máximo de pose & \SI{0.045}{m}\\
Distancia mínima positiva a obstáculo & \SI{0.056}{m}\\
Riesgo máximo de obstáculo & 0,829\\
\bottomrule
\end{longtable}

Junto con las métricas, se comprobó la ejecución en la escena: Bill se mueve
entre WS1 y WS2, toma y entrega piezas, trabaja en la mesa, aparecen
estanterías, mesas y sofá, T1 y T2 tienen geometría reconocible, la cola de
tareas queda registrada, el robot vuelve a carga y el planificador activa
\texttt{SLAM\_WAYPOINT} cuando detecta un bloqueo local.

\section{Desempeño del estimador Kalman-landmark}

En la exportación específica del estimador, la escena usa PID como controlador
base y \texttt{KALMAN\_LANDMARK} como algoritmo SLAM. Los resultados fueron:

\begin{itemize}
  \item RMSE de posición: \SI{0.03149}{m}.
  \item Error medio de posición: \SI{0.02981}{m}.
  \item Error P95: \SI{0.04760}{m}.
  \item Error máximo: \SI{0.05626}{m}.
  \item RMSE de orientación: \SI{0.01987}{rad}.
  \item Covarianza media estimada: 0,00081.
  \item Landmarks finales: 20.
  \item Actualizaciones Kalman/SLAM: 4295.
\end{itemize}

El número de elementos del mapa depende de la representación. En
Kalman-landmark son rasgos puntuales fusionados; en las variantes de grid son
celdas ocupadas o grupos de celdas. La validación funcional cerró con 24 rasgos
y la exportación específica Kalman-landmark con 20, mientras que las grids
acabaron con mapas más densos.

\begin{figure}[H]
  \centering
  \includegraphics[width=\textwidth]{phase1_slam_estimator_performance.png}
  \caption{Resultados del estimador Kalman-landmark.}
  \label{fig:estimator-performance}
\end{figure}

\section{Comparación de algoritmos SLAM}

La comparación usa el mismo escenario, sensores, tareas y controlador para los
cuatro modos de SLAM. La Tabla~\ref{tab:slam-comparison} recoge los
valores principales.

Los tiempos y longitudes de ruta salen parecidos porque el plan de misión, los
puntos de entrega y el controlador de movimiento son los mismos. Cambia la
representación del mapa usada por los bloqueos y los puntos intermedios. Por
eso la columna de rasgos exige cuidado: en Kalman son landmarks fusionadas; en
las grids son celdas o grupos de celdas ocupadas.

\begingroup
\footnotesize
\setlength{\tabcolsep}{3pt}
\begin{longtable}{P{0.26\textwidth}P{0.16\textwidth}P{0.16\textwidth}P{0.18\textwidth}P{0.16\textwidth}}
\caption{Comparación SLAM en CoppeliaSim.}\label{tab:slam-comparison}\\
\toprule
\textbf{Métrica} & \textbf{GMapping} & \textbf{Hector} & \textbf{Cartographer} & \textbf{Kalman}\\
\midrule
\endfirsthead
\toprule
\textbf{Métrica} & \textbf{GMapping} & \textbf{Hector} & \textbf{Cartographer} & \textbf{Kalman}\\
\midrule
\endhead
Completado & Sí & Sí & Sí & Sí\\
Duración [s] & 125,60 & 123,55 & 125,25 & \textbf{121,45}\\
Ruta [m] & 15,997 & \textbf{15,233} & 15,968 & 15,763\\
RMSE pos. [m] & 0,04117 & 0,04445 & 0,04149 & \textbf{0,03142}\\
Error P95 [m] & 0,05718 & 0,07596 & 0,06026 & \textbf{0,04670}\\
Error max. [m] & 0,07354 & 0,10862 & 0,07362 & \textbf{0,05824}\\
Rasgos finales & 94 & 84 & 91 & 20\\
Actualizaciones SLAM & 4514 & 4470 & 4474 & 4275\\
Submapas & 0 & 0 & \textbf{13} & 0\\
Cierres de ciclo & 0 & 0 & \textbf{4} & 0\\
Batería usada [\%] & 24,62 & \textbf{23,63} & 24,50 & 24,02\\
\bottomrule
\end{longtable}
\endgroup

\begin{figure}[H]
  \centering
  \includegraphics[width=\textwidth]{phase1_slam_algorithm_comparison.png}
  \caption{Comparación gráfica de los cuatro modos SLAM: trayectoria, precisión, crecimiento del mapa y resumen de error.}
  \label{fig:slam-comparison}
\end{figure}

\section{Interpretación de resultados}

Kalman-landmark obtuvo el menor error de pose y el menor consumo. En una celda
pequeña, las landmarks fusionadas son suficientes para navegar. GMapping y Hector
generan mapas más densos y comparan grid de ocupación con
\textit{scan matching}; pagan con algo más de error y más trabajo de ajuste.
Cartographer-submap redujo la ruta frente a GMapping y Hector y registró
estructura por submapas con 4 cierres de ciclo locales, aunque su RMSE quedó por
encima de Kalman-landmark.

Para la escena base usamos \texttt{KALMAN\_LANDMARK} como estimador base. Los otros
tres métodos quedan para comparar mapas calculados con los mismos rayos de
proximidad.
